\documentclass[UTF8,12pt,openany]{ctexrep}

\usepackage[a4paper,top=2.35cm,bottom=2.3cm,left=2.35cm,right=2.35cm,headheight=16pt]{geometry}
\usepackage{fontspec}
\setmainfont{TeX Gyre Pagella}
\setsansfont{TeX Gyre Heros}
\setmonofont{DejaVu Sans Mono}[Scale=0.84]
\setCJKmainfont{Noto Serif CJK SC}
\setCJKsansfont{Noto Sans CJK SC}
\setCJKmonofont{Noto Sans Mono CJK SC}

\usepackage{amsmath,amssymb,bm}
\usepackage{graphicx}
\usepackage{booktabs,tabularx,longtable,array,multirow}
\usepackage{siunitx}
\usepackage[dvipsnames]{xcolor}
\usepackage{tikz}
\usepackage[most]{tcolorbox}
\usepackage{enumitem}
\usepackage{caption,subcaption}
\usepackage{float}
\usepackage{fancyhdr}
\usepackage{titlesec}
\usepackage{listings}
\usepackage{xurl}
\usepackage{hyperref}
\usepackage{lastpage}

\definecolor{Navy}{HTML}{153F66}
\definecolor{Blue}{HTML}{2F6F9F}
\definecolor{Teal}{HTML}{2A8C82}
\definecolor{Green}{HTML}{4D8F62}
\definecolor{Orange}{HTML}{D9822B}
\definecolor{Red}{HTML}{B04A4A}
\definecolor{Purple}{HTML}{745AA3}
\definecolor{LightBlue}{HTML}{EAF2F8}
\definecolor{LightGreen}{HTML}{EAF4ED}
\definecolor{LightOrange}{HTML}{FBF0E4}
\definecolor{LightGray}{HTML}{F3F5F7}
\definecolor{Dark}{HTML}{24313D}

\hypersetup{
  colorlinks=true,
  linkcolor=Navy,
  urlcolor=Blue,
  citecolor=Teal,
  pdftitle={1 kHz 深海声传播损失快速代理模型技术白皮书},
  pdfauthor={Ocean Acoustic Surrogate Project},
  pdfsubject={Bellhop 数据生成、FNO 代理建模与验收报告}
}

\graphicspath{{assets/}}
\captionsetup{font=small,labelfont={bf,color=Navy}}
\setlist{leftmargin=2em,itemsep=0.25em,topsep=0.35em}
\renewcommand{\arraystretch}{1.28}
\setlength{\tabcolsep}{5pt}
\setcounter{tocdepth}{2}
\setcounter{secnumdepth}{3}

\titleformat{\chapter}[display]
  {\bfseries\sffamily\color{Navy}}
  {\fontsize{16}{18}\selectfont 第\,\thechapter\,章}
  {0.4em}
  {\fontsize{25}{30}\selectfont}
  [\vspace{0.35em}\color{Blue}\titlerule]
\titleformat{\section}{\Large\bfseries\sffamily\color{Navy}}{\thesection}{0.75em}{}
\titleformat{\subsection}{\large\bfseries\sffamily\color{Blue}}{\thesubsection}{0.65em}{}

\pagestyle{fancy}
\fancyhf{}
\fancyhead[L]{\small\sffamily\color{Navy}1 kHz 深海声传播损失快速代理模型}
\fancyhead[R]{\small\sffamily\color{Navy}技术白皮书 V1.0}
\fancyfoot[L]{\small\color{gray}Ocean Acoustic Surrogate Project}
\fancyfoot[R]{\small\color{gray}第 \thepage\ 页 / 共 \pageref{LastPage} 页}
\renewcommand{\headrulewidth}{0.5pt}
\renewcommand{\footrulewidth}{0pt}

\newcolumntype{Y}{>{\raggedright\arraybackslash}X}
\newcolumntype{C}[1]{>{\centering\arraybackslash}p{#1}}
\newcolumntype{L}[1]{>{\raggedright\arraybackslash}p{#1}}

\newtcolorbox{executivebox}[1][]{
  enhanced,
  breakable,
  colback=LightBlue,
  colframe=Blue,
  boxrule=0.8pt,
  arc=1.5mm,
  left=4mm,right=4mm,top=3mm,bottom=3mm,
  title=#1,
  fonttitle=\bfseries\sffamily
}
\newtcolorbox{passbox}[1][]{
  enhanced,
  breakable,
  colback=LightGreen,
  colframe=Green,
  boxrule=0.8pt,
  arc=1.5mm,
  left=4mm,right=4mm,top=3mm,bottom=3mm,
  title=#1,
  fonttitle=\bfseries\sffamily
}
\newtcolorbox{riskbox}[1][]{
  enhanced,
  breakable,
  colback=LightOrange,
  colframe=Orange,
  boxrule=0.8pt,
  arc=1.5mm,
  left=4mm,right=4mm,top=3mm,bottom=3mm,
  title=#1,
  fonttitle=\bfseries\sffamily
}

\lstdefinestyle{command}{
  basicstyle=\ttfamily\small,
  backgroundcolor=\color{LightGray},
  frame=single,
  rulecolor=\color{gray!35},
  breaklines=true,
  columns=fullflexible,
  xleftmargin=1em,
  framexleftmargin=0.7em,
  aboveskip=0.7em,
  belowskip=0.7em
}

\newcommand{\metric}[2]{\textbf{#1}\;#2}
\newcommand{\code}[1]{\texttt{#1}}
\newcommand{\figfull}[3]{
  \begin{figure}[H]
    \centering
    \includegraphics[width=#1\textwidth]{#2}
    \caption{#3}
  \end{figure}
}

\begin{document}

% -----------------------------------------------------------------------------
% Cover
% -----------------------------------------------------------------------------
\begin{titlepage}
\begin{tikzpicture}[remember picture,overlay]
  \fill[Navy] (current page.north west) rectangle ([yshift=-10.7cm]current page.north east);
  \fill[Blue] ([yshift=-10.7cm]current page.north west) rectangle ([yshift=-10.88cm]current page.north east);
  \fill[LightBlue] ([yshift=3.2cm]current page.south west) rectangle (current page.south east);
\end{tikzpicture}

\vspace*{1.25cm}
{\sffamily\color{white}\Large 项目技术白皮书}\par
\vspace{1.2cm}
{\sffamily\bfseries\color{white}\fontsize{30}{38}\selectfont
1 kHz 深海声传播损失\par
快速代理模型}\par
\vspace{0.55cm}
{\sffamily\color{white}\Large
基于高质量 Bellhop 标签与各向异性 FNO}\par

\vfill

\begin{tcolorbox}[colback=white,colframe=Navy,boxrule=1pt,arc=2mm,width=0.78\textwidth]
\centering
\begin{tabular}{L{3.5cm}L{6.8cm}}
\textbf{文档版本} & V1.0 / MVP 验收版 \\
\textbf{实验日期} & 2026-08-27 \\
\textbf{数据集} & \code{scs\_june\_narrow\_v0.1} \\
\textbf{模型版本} & \code{mvp-v0.1} \\
\textbf{交付状态} & \textcolor{Green}{\textbf{已达到 2 dB / 100 ms 核心指标}} \\
\end{tabular}
\end{tcolorbox}

\vfill
{\sffamily\color{Navy}\large Ocean Acoustic Surrogate Project}\par
\vspace{0.25cm}
{\color{Dark}面向固定 1 kHz、50 km、典型深海环境族的首阶段可行性与交付验证}\par
\vspace{1.25cm}
\end{titlepage}

\pagenumbering{roman}

% -----------------------------------------------------------------------------
% Document control
% -----------------------------------------------------------------------------
\chapter*{文档说明}
\addcontentsline{toc}{chapter}{文档说明}

本白皮书面向项目甲方、技术评审专家和后续工程实施人员。它回答五个核心问题：

\begin{enumerate}
  \item 数值标签如何生成，为什么可以认为标签质量足够高；
  \item 声速剖面如何变成模型输入，模型为什么采用 Fourier Neural Operator；
  \item 如何避免数据泄漏，如何定义 baseline、训练、验证和密封测试；
  \item 五轮消融分别验证了什么，哪些复杂设计没有产生净收益；
  \item 最终模型是否在独立重载后仍满足精度和时延指标，以及成果适用边界是什么。
\end{enumerate}

\begin{riskbox}[重要口径]
本阶段的目标是\textbf{先在一个简单、固定、严格满足硬参数的环境族中达标}，不以跨海区、
跨频率、跨源深或跨底质复用为第一目标。因此，本文所有“通过”结论均限定在第
\ref{chap:scope}章给出的冻结验收域内。甲方自有 Bellhop 到位后，仍需先完成同 case 数值对齐。
\end{riskbox}

\section*{版本记录}
\noindent
\begin{tabularx}{\textwidth}{L{2.3cm}L{2.5cm}Y Y}
\toprule
版本 & 日期 & 内容 & 状态 \\
\midrule
V0.1 & 2026-08-27 & 数据契约、收敛审计、512 样本、五轮实验 & 已封存 \\
V1.0 & 2026-08-27 & 甲方友好白皮书、完整公式、消融和独立验收 & 本次交付 \\
\bottomrule
\end{tabularx}

\clearpage
\tableofcontents
\clearpage
\pagenumbering{arabic}

% -----------------------------------------------------------------------------
\chapter{执行摘要}
\label{chap:executive}

\begin{executivebox}[结论先行]
本项目已建立一条从\textbf{窄域深海 SSP 采样、高质量 Bellhop 非相干 TL 生成、数据封存、
FNO 训练、密封测试到独立新进程验收}的完整闭环。正式数据共 512 场、25,600 rays/场、
失败数为 0。精度优胜模型在 A100 上达到测试 RMSE 0.5165 dB、MAE 0.1802 dB、P95
端到端推理 4.73 ms；便携小模型在独立 CPU 重算中达到 RMSE 0.7582 dB、MAE 0.3067 dB、
P95 38.32 ms。两套模型均显著优于 2 dB / 100 ms 硬门槛。
\end{executivebox}

\section{项目问题与最终答案}

项目要求在频率 1 kHz、距离 50 km 的典型非均匀深海环境中，以神经算子快速预测二维
传播损失场。传统 Bellhop 每场标签的中位求解时间约 9.14 s；代理模型在 A100 上一次完整
推理约 4.73 ms，观测墙钟比约为 1,930 倍。更重要的是，代理预测不是只在平均意义上“看起来
相似”：64 条密封测试样本中最差逐场 RMSE 仍只有 0.769 dB。

\begin{passbox}[核心验收证书]
\begin{tabularx}{\textwidth}{Y C{2.5cm} C{2.5cm} C{2.2cm}}
\toprule
验收项 & 上限 & 实测 & 判定 \\
\midrule
R3 / CUDA 测试 RMSE & 2.000 dB & 0.5165 dB & \textcolor{Green}{\textbf{通过}} \\
R3 / CUDA 测试 MAE & 2.000 dB & 0.1802 dB & \textcolor{Green}{\textbf{通过}} \\
R3 / CUDA 完整推理 P95 & 100 ms & 4.73 ms & \textcolor{Green}{\textbf{通过}} \\
R1 / CPU 独立测试 RMSE & 2.000 dB & 0.7582 dB & \textcolor{Green}{\textbf{通过}} \\
R1 / CPU 独立完整推理 P95 & 100 ms & 38.32 ms & \textcolor{Green}{\textbf{通过}} \\
\bottomrule
\end{tabularx}
\end{passbox}

\figfull{0.98}{fig10_acceptance_margin.pdf}{双模型相对验收门槛的余量。每根柱表示实测值占验收上限的比例。}

\section{推荐交付策略}

最终保留两个模型不是重复建设，而是针对不同部署条件提供明确选择：

\begin{itemize}
  \item \textbf{R3 精度优胜模型：}6.30 M 参数，适用于已有 CUDA GPU 的甲方环境；测试
  RMSE 最低，A100 P95 约 4--5 ms。
  \item \textbf{R1 便携 MVP：}1.33 M 参数，checkpoint 约 11 MB，在 12 线程 CPU 上也可
  低于 100 ms；适用于硬件条件尚未冻结或需要最稳妥首版集成的场景。
\end{itemize}

\figfull{0.98}{fig01_project_pipeline.pdf}{项目从环境采样到独立验收的端到端闭环。}

% -----------------------------------------------------------------------------
\chapter{需求拆解与验收边界}
\label{chap:scope}

\section{硬约束}

本阶段把需求拆成环境、输出、数值参考和在线性能四类约束。

\begin{table}[H]
\centering
\caption{冻结验收域}
\begin{tabularx}{\textwidth}{L{3.5cm}Y L{3.4cm}}
\toprule
类别 & 冻结设置 & 备注 \\
\midrule
频率 & 1000 Hz & 单频，不做跨频率泛化 \\
传播距离 & 0.1--50 km，256 点 & 最大距离严格覆盖 50 km \\
水深 & 2000 m & 接收深度 10--1990 m，96 点 \\
声源深度 & 50 m & 不做深源/浅源迁移 \\
海洋环境 & 距离无关、随深度非均匀 SSP & 典型深海声道 \\
海底 & 流体沙质半空间 & 1700 m/s，2000 kg/m\textsuperscript{3}，0.8 dB/$\lambda$ \\
参考解 & Bellhop 非相干 TL & 25,600 rays \\
精度 & 测试 RMSE、MAE 均不超过 2 dB & 只在有效网格计分 \\
时延 & batch=1 热态完整推理 P95 不超过 100 ms & 明确报告硬件和计时边界 \\
\bottomrule
\end{tabularx}
\end{table}

\section{主动不扩展的维度}

为了避免“数据覆盖越多、任务反而越难”的风险，首阶段没有加入下列变化：不同源深、不同
频率、不同水深、不同底质、距离相关 SSP、海底地形、海面状态、跨季节海区或相干/半相干
场。该策略不代表这些维度不重要，而是把它们从\textbf{首个可验收闭环}中剥离。

\section{成功标准的工程化解释}

“平均误差不超过 2 dB”可能被理解为 MAE，也可能被理解为 RMSE。为避免口径争议，项目
同时要求二者均通过，并额外报告 bias、P95 绝对误差、逐样本 RMSE、最差样本 RMSE 和
高梯度区域 RMSE。“模型计算时间”则采用包含特征 tensor 创建、设备传输、模型前向、反归一化、
复制回 CPU NumPy 和 GPU 同步的端到端计时；不把模型加载和首次 JIT/库初始化混入热态延迟。

% -----------------------------------------------------------------------------
\chapter{物理量、数值参考与标签定义}
\label{chap:physics}

\section{传播损失定义}

设接收位置为 $\bm{x}=(z,r)$，复声压为 $p(\bm{x})$，参考声压为 $p_{\mathrm{ref}}$。相干
传播损失通常写为
\begin{equation}
  TL(\bm{x})=-20\log_{10}\left(\frac{|p(\bm{x})|}{|p_{\mathrm{ref}}|}\right)\quad\text{dB}.
\end{equation}

本项目采用 Bellhop 的\textbf{非相干}场输出。对于多条到达路径，不保留路径间相位干涉，
而在强度层面累加：
\begin{equation}
  I_{\mathrm{inc}}(\bm{x})=\sum_{j=1}^{N_a}|p_j(\bm{x})|^2,
  \qquad
  TL_{\mathrm{inc}}(\bm{x})=-10\log_{10}\left(\frac{I_{\mathrm{inc}}(\bm{x})}{I_{\mathrm{ref}}}\right).
\end{equation}

这种定义与学生前期实验保持一致，同时比半相干场对离散射线数更稳定。本文所有标签、训练
目标和指标都明确指向非相干 TL，不混用相干或半相干结果。

\section{Bellhop 数值环境}

Bellhop 通过射线追踪和高斯束近似在给定 SSP、源深、边界和接收网格上计算声场。每个正式
样本包含：输入 SSP、Bellhop case、二维 TL、距离/深度坐标、有效掩膜、环境参数、求解
耗时、数据 split、配置哈希和依赖 Git SHA。这样可以从任意数据样本追溯到数值 case，而不
只是保留一个无法解释的训练数组。

\begin{riskbox}[为什么必须做射线数收敛]
“使用 Bellhop”并不自动意味着标签足够准确。射线数、射线角、输出模式和网格都会影响
离散结果。如果标签自身在 1--2 dB 量级波动，代理模型再低的训练 loss 也没有意义。因此，
正式生成前先把射线数作为唯一变量进行收敛审计。
\end{riskbox}

\section{射线数收敛审计}

在 8 个独立 SSP 上分别运行 3,200、6,400、12,800 和 25,600 rays。定义相邻两级标签差
的逐场 RMSE 为
\begin{equation}
  E_{n\rightarrow 2n}^{(s)}=\sqrt{\frac{1}{|\Omega_s|}
  \sum_{(z,r)\in\Omega_s}\left(TL_n^{(s)}(z,r)-TL_{2n}^{(s)}(z,r)\right)^2}.
\end{equation}

\begin{table}[H]
\centering
\caption{Bellhop 射线数收敛结果}
\begin{tabular}{C{4.0cm}C{3.6cm}C{3.6cm}}
\toprule
比较 & 8 个 SSP 平均 RMSE & 最差逐场 RMSE \\
\midrule
3,200 $\rightarrow$ 6,400 & 0.1039 dB & 0.1484 dB \\
6,400 $\rightarrow$ 12,800 & 0.0656 dB & 0.1239 dB \\
12,800 $\rightarrow$ 25,600 & 0.0520 dB & 0.1410 dB \\
\bottomrule
\end{tabular}
\end{table}

最高两级的差异比最终 2 dB 门槛低一个数量级以上，因此把 25,600 rays 固定为正式标签
设置。该结论仅对当前非相干 TL、环境族和网格成立。

\figfull{0.98}{fig03_label_quality.pdf}{标签收敛误差与 512 场正式数据的生成稳定性。}

% -----------------------------------------------------------------------------
\chapter{数据集生成与质量控制}
\label{chap:data}

\section{基础 SSP}

基础声速剖面由 11 个深度控制点定义：
\begin{center}
\small
\begin{tabular}{r|rrrrrrrrrrr}
深度 (m) & 0 & 50 & 100 & 200 & 400 & 700 & 1000 & 1300 & 1600 & 1800 & 2000 \\
\hline
声速 (m/s) & 1545 & 1537 & 1529 & 1516 & 1501 & 1491 & 1487 & 1488 & 1491 & 1493 & 1496 \\
\end{tabular}
\end{center}

先采用自然三次样条得到平滑 $c_0(z)$，再用四个低维参数产生物理上可解释的变化：
\begin{align}
c(z;\bm{\theta})=
&\;c_0\!\left(\operatorname{clip}(z-\Delta z,0,2000)\right)+\Delta c \\
&+ A_t\exp\left[-\frac{1}{2}\left(\frac{z-250}{220}\right)^2\right]
+G_d\operatorname{clip}\left(\frac{z-1000}{1000},0,1\right),
\end{align}
其中 $\bm{\theta}=(\Delta c,A_t,\Delta z,G_d)$。

\begin{table}[H]
\centering
\caption{SSP 参数范围}
\begin{tabularx}{\textwidth}{L{4.0cm}L{3.2cm}L{3.4cm}Y}
\toprule
参数 & 冻结范围 & 正式样本实际范围 & 物理作用 \\
\midrule
全局偏移 $\Delta c$ & $[-1.5,1.5]$ m/s & $[-1.4944,1.4953]$ & 整体声速平移 \\
温跃层幅度 $A_t$ & $[-2.0,2.0]$ m/s & $[-1.9980,1.9952]$ & 改变上层梯度 \\
声道轴移动 $\Delta z$ & $[-100,100]$ m & $[-99.82,99.82]$ & 垂向移动声道结构 \\
深层梯度 $G_d$ & $[-1.5,1.5]$ m/s & $[-1.4965,1.4945]$ & 改变 1000 m 以下趋势 \\
\bottomrule
\end{tabularx}
\end{table}

所有生成声速被限制在 1475--1565 m/s 内，不叠加逐深度白噪声，从而避免产生不必要的
高频非物理结构。

\section{Latin hypercube 采样}

对四维单位超立方体使用 seed 20260827 的优化 Latin hypercube。设第 $i$ 个单位样本为
$\bm{u}_i\in[0,1]^4$，第 $j$ 个参数上下界为 $[a_j,b_j]$，则
\begin{equation}
  \theta_{ij}=a_j+u_{ij}(b_j-a_j).
\end{equation}
Latin hypercube 保证每个单参数边缘区间近似均匀覆盖，同时只需 512 个样本即可获得比纯
随机采样更稳定的空间填充。它的目的不是覆盖所有海洋，而是高效覆盖本次冻结小族。

\figfull{0.98}{fig02_dataset_design.pdf}{512 条 SSP 的剖面、四参数边缘分布和固定数据划分。}

\section{逐样本生成流程}

每个 $\bm{\theta}_i$ 的标签生成步骤如下：
\begin{enumerate}
  \item 由四参数构造 33 点平滑 SSP；
  \item 写入固定水深、源深、频率、海底和接收网格；
  \item 调用 Bellhop，明确指定 \code{field\_mode=incoherent} 和 25,600 rays；
  \item 读取二维 TL、距离、深度和 case 路径；
  \item 对非有限位置生成 \code{valid\_mask}，数值数组使用 120 dB 填充但不计分；
  \item 保存 \code{sample.npz} 和 \code{metadata.json}，支持逐样本断点恢复；
  \item 全部成功后按固定顺序打包，并计算 SHA-256 和 manifest。
\end{enumerate}

正式生成途中共享盘曾返回用户配额错误。已完成的 36 条样本被整体移动到本地大容量目录，
随后利用逐样本缓存断点续跑；没有删除、重算或改变参数。最终 512/512 成功。

\section{划分策略与数据泄漏控制}

对 512 个已采样 SSP 做 seed 固定的确定性随机排列，得到 384/64/64 的训练、验证、测试
划分。三者同分布但样本互斥。训练集用于参数优化和目标变换拟合；验证集用于每个 epoch 的
checkpoint 选择；测试集只用于冻结配置的最终评估。五轮消融复用同一划分，避免通过更换
测试样本获得虚假提升。

\section{正式数据统计}

\begin{table}[H]
\centering
\caption{正式数据集统计}
\begin{tabularx}{\textwidth}{L{5.1cm}Y}
\toprule
项目 & 数值 \\
\midrule
样本数 / 失败数 & 512 / 0 \\
单场输出 & 96 深度 $\times$ 256 距离 = 24,576 网格 \\
有效网格覆盖率 & 99.6877\% \\
TL 有效值范围 & 37.969--94.137 dB \\
TL 有效值均值 / 标准差 & 81.335 / 8.253 dB \\
Bellhop 累计时间 & 4,724.25 s（78.74 min） \\
单场 Bellhop 中位 / P95 & 9.138 / 9.492 s \\
压缩数据集大小 & 41,170,478 bytes \\
含完整 case 的目录 & 约 796 MB \\
数据 SHA-256 & \footnotesize\nolinkurl{13f4246b86a2fff33f2e5d587c1443be47eb15b4ad3d5c413688464ed8c7281d} \\
\bottomrule
\end{tabularx}
\end{table}

% -----------------------------------------------------------------------------
\chapter{数据处理与特征工程}
\label{chap:processing}

\section{SSP 插值与网格广播}

网络输出定义在 96$\times$256 接收网格，而 SSP 是一维函数。首先在接收深度 $z_m$ 上
线性插值得到 $c_i(z_m)$，再沿距离方向广播：
\begin{equation}
  X^{\mathrm{ssp}}_i(z_m,r_n)=\frac{c_i(z_m)-1500}{50}.
\end{equation}
该归一化把输入控制在约 $[-0.5,1]$ 的稳定范围，减少优化器需要处理的绝对量级差异。

\section{位置编码}

FNO 内部为每个网格增加归一化深度和距离坐标：
\begin{equation}
  X^z(z_m,r_n)=\frac{m}{N_z-1},\qquad
  X^r(z_m,r_n)=\frac{n}{N_r-1}.
\end{equation}
坐标使网络能够区分同一 SSP 特征在不同空间位置的物理含义。

\section{Hankel 附加特征}

R4 和 R5 参考学生设置，加入二维圆柱 Green 函数的 Hankel 幅度：
\begin{equation}
  h(r)=\operatorname{Normalize}\left[\log_{10}\left|H_0^{(1)}(kr)\right|\right],
  \qquad k=\frac{2\pi f}{1500}.
\end{equation}
它提供近似几何扩散先验，但消融表明在本窄域中并未改善整体 RMSE，因此最终 R3 不使用
该通道。

\section{训练均值场残差变换}

直接预测 40--100 dB 的完整 TL 会让网络大量容量用于学习几乎固定的传播趋势。为此，先只
用训练集有效网格拟合逐位置平均场：
\begin{equation}
  \overline{TL}_{\mathrm{train}}(z,r)=
  \frac{\sum_{i\in\mathcal{D}_{\mathrm{train}}}M_i(z,r)TL_i(z,r)}
  {\sum_{i\in\mathcal{D}_{\mathrm{train}}}M_i(z,r)}.
\end{equation}
随后定义残差尺度 $\sigma_{\mathrm{res}}$，训练目标为
\begin{equation}
  Y_i(z,r)=\frac{TL_i(z,r)-\overline{TL}_{\mathrm{train}}(z,r)}
  {\max(\sigma_{\mathrm{res}},1\ \mathrm{dB})}.
\end{equation}
推理时执行逆变换：
\begin{equation}
  \widehat{TL}_i(z,r)=\overline{TL}_{\mathrm{train}}(z,r)
  +\sigma_{\mathrm{res}}\widehat{Y}_i(z,r).
\end{equation}
均值场和尺度都不读取验证或测试标签，从根源上避免目标统计泄漏。

\section{有效掩膜}

对 Bellhop 非有限或不可计分的位置定义 $M_i(z,r)=0$；其他位置为 1。所有 loss 和指标
都以 $M$ 为准。120 dB 只是数组填充值，不参与训练或验收。这一点非常重要：把填充值当成
真实 TL 会人为制造巨大误差和错误梯度。

% -----------------------------------------------------------------------------
\chapter{代理模型方法}
\label{chap:model}

\section{算子学习问题}

目标不是一个有限维分类或回归，而是学习函数到函数的映射：
\begin{equation}
  \mathcal{G}_{\theta}:c(z)\longmapsto TL(z,r),
\end{equation}
其中输入是声速函数，输出是二维声场函数。Fourier Neural Operator 在频域参数化全局积分
算子，能够用少量层捕获远距离空间关联，适合本问题中的大尺度传播结构。

\section{FNO 谱卷积}

令第 $\ell$ 层特征为 $v_\ell\in\mathbb{R}^{C\times N_z\times N_r}$。二维离散 Fourier
变换为 $\mathcal{F}$，截断模态集合为 $\mathcal{K}$，复权重为
$R_\ell(k_z,k_r)\in\mathbb{C}^{C\times C}$，则谱分支为
\begin{equation}
  \mathcal{K}_{\phi_\ell}(v_\ell)=
  \mathcal{F}^{-1}\left(
  \mathbf{1}_{(k_z,k_r)\in\mathcal{K}}
  R_\ell(k_z,k_r)\,\mathcal{F}(v_\ell)(k_z,k_r)
  \right).
\end{equation}
深度和距离方向物理尺度不同，因此使用各向异性截断 $K_z\neq K_r$。

\section{单个 FNO block}

基础 block 同时保留频域全局分支和空间局部分支：
\begin{align}
  u_\ell &= \mathcal{K}_{\phi_\ell}(v_\ell)+W_\ell v_\ell,\\
  v_{\ell+1} &= \operatorname{GELU}\left(
  \operatorname{GroupNorm}(u_\ell)+\rho_\ell v_\ell\right),
\end{align}
其中 $W_\ell$ 是 1$\times$1 卷积；若启用局部增强，则替换为 depthwise 3$\times$3
卷积和点卷积；$\rho_\ell\in\{0,1\}$ 控制 block 残差。

\section{非周期 padding}

普通 FFT 隐含周期边界假设，而距离和深度方向都不是周期场。R3 在深度尾部复制 padding 8
点、距离尾部复制 padding 16 点，经过四层后裁剪回原网格，以减少边界环绕伪影。

\figfull{0.98}{fig04_model_architecture.pdf}{R3 各向异性 FNO 的输入、谱层和残差解码结构。}

\section{两套最终模型}

\begin{table}[H]
\centering
\caption{最终保留模型配置}
\begin{tabularx}{\textwidth}{L{4.2cm}C{3.6cm}C{3.6cm}Y}
\toprule
配置 & R1 便携 MVP & R3 精度优胜 & 说明 \\
\midrule
hidden channels & 24 & 32 & 内部特征宽度 \\
深度/距离 modes & 12 / 24 & 16 / 48 & 各向异性截断 \\
FNO 层数 & 4 & 4 & 保持一致 \\
padding 深度/距离 & 0 / 0 & 8 / 16 & R3 抑制非周期边界 \\
block 残差 & 否 & 是 & R3 启用 \\
Hankel 特征 & 否 & 否 & 消融后不保留 \\
参数量 & 1,333,121 & 6,300,417 & R3 约 4.7 倍 \\
checkpoint & 约 11 MB & 约 49 MB & FP32 \\
\bottomrule
\end{tabularx}
\end{table}

% -----------------------------------------------------------------------------
\chapter{训练、损失函数与模型选择}
\label{chap:training}

\section{基础值损失}

设归一化预测为 $\widehat{Y}$，目标为 $Y$，有效掩膜为 $M$，基础 MSE 为
\begin{equation}
  \mathcal{L}_{\mathrm{value}}=
  \frac{\sum_{i,z,r}M_i(z,r)\left(\widehat{Y}_i(z,r)-Y_i(z,r)\right)^2}
  {\sum_{i,z,r}M_i(z,r)}.
\end{equation}

\section{梯度损失}

为尝试强化会聚边缘，R5 增加深度和距离的一阶差分匹配：
\begin{equation}
  \mathcal{L}_{\mathrm{grad}}=\frac{1}{2}
  \left[
  \operatorname{MSE}_M(\Delta_z\widehat{Y},\Delta_zY)+
  \operatorname{MSE}_M(\Delta_r\widehat{Y},\Delta_rY)
  \right].
\end{equation}

\section{高梯度网格加权}

由参考目标定义局部结构强度
\begin{equation}
  g_i(z,r)=|\Delta_zY_i(z,r)|+|\Delta_rY_i(z,r)|.
\end{equation}
在每个 batch 内取有效 $g$ 的 90\% 分位数 $q_{0.9}$，定义困难网格
$D_i=M_i\mathbf{1}(g_i\ge q_{0.9})$，其附加损失为
\begin{equation}
  \mathcal{L}_{\mathrm{high}}=
  \frac{\sum D_i(\widehat{Y}_i-Y_i)^2}{\sum D_i}.
\end{equation}
R5 总目标为
\begin{equation}
  \mathcal{L}=\mathcal{L}_{\mathrm{value}}+0.05\mathcal{L}_{\mathrm{grad}}
  +1.0\mathcal{L}_{\mathrm{high}}.
\end{equation}
由于 R5 多了两个非负项，其 loss 数值不能与 R1--R4 直接横向比较；最终必须回到相同 dB
指标判断是否有净收益。

\section{统一训练预算}

\begin{table}[H]
\centering
\caption{训练超参数}
\begin{tabularx}{\textwidth}{L{4.4cm}Y L{4.2cm}}
\toprule
项目 & 设置 & 目的 \\
\midrule
随机种子 & 20260827 & 可复现初始化与 shuffle \\
优化器 & AdamW & 稳定训练谱权重 \\
初始学习率 & $10^{-3}$ & 五轮统一 \\
weight decay & $10^{-5}$ & 轻量正则 \\
调度 & cosine，最低为初始的 1/50 & 后期精细收敛 \\
batch size & 16 & A100 显存与效率平衡 \\
最大 epoch & 180 & 小于学生 2000 epochs \\
early stopping patience & 35 & 防止无效延长 \\
梯度裁剪 & 1.0 & 抑制偶发谱层尖峰 \\
模型选择 & 最低 validation RMSE & 测试集不参与 checkpoint 选择 \\
\bottomrule
\end{tabularx}
\end{table}

\figfull{0.98}{fig05_training_curves.pdf}{五轮训练 loss 与验证 RMSE。R5 loss 因额外项而处于不同量级。}

% -----------------------------------------------------------------------------
\chapter{Baseline 与测试协议}
\label{chap:test}

\section{训练均值场 baseline}

最简单 baseline 是忽略输入 SSP，对所有样本输出训练集逐网格平均 TL。它没有可训练神经
网络参数，但完整遵守“测试标签不可用于统计”的规则。其测试结果为：RMSE 1.4211 dB、MAE
0.8046 dB、最差单样本 RMSE 1.9101 dB。

该 baseline 本身已经低于 2 dB。这不是应被隐藏的问题，而是窄域数据设计成功的证据。
代理模型需要回答的更强问题是：\textbf{能否利用 SSP 把误差从 1.42 dB 进一步压到约 0.52
dB，并准确移动会聚结构。}

\section{学生方案参考 baseline}

学生前期设置使用 3,060 条南海东侧 6 月 SSP、64 通道、72 modes、4 个 Fourier 层、
Hankel 幅度特征和 2000 epochs，随机 9:1 测试 RMSE 约 1.1--1.2 dB。它提供了技术可达性
参考，但由于数据、划分和实现不同，不作为严格同表数值 baseline。本项目重点验证能否以
更少数据、更少 epoch 和更清晰的标签契约获得更大验收余量。

\section{主要指标公式}

对测试有效网格集合 $\Omega$，误差 $e_j=\widehat{TL}_j-TL_j$：
\begin{align}
  \mathrm{RMSE} &= \sqrt{\frac{1}{|\Omega|}\sum_{j\in\Omega}e_j^2},\\
  \mathrm{MAE} &= \frac{1}{|\Omega|}\sum_{j\in\Omega}|e_j|,\\
  \mathrm{Bias} &= \frac{1}{|\Omega|}\sum_{j\in\Omega}e_j,\\
  P_q &= \operatorname{Quantile}_q\left(\{|e_j|:j\in\Omega\}\right).
\end{align}

整体 RMSE 是将全部样本有效网格合并后的 micro 指标。对第 $s$ 条样本另计算
\begin{equation}
  \mathrm{RMSE}_s=\sqrt{\frac{1}{|\Omega_s|}\sum_{j\in\Omega_s}e_{sj}^2},
\end{equation}
再报告中位、P90 和最大值，以避免大量简单网格掩盖困难样本。

\section{高梯度区域指标}

在测试参考 TL 上计算
\begin{equation}
  g=|\Delta_zTL|+|\Delta_rTL|,
\end{equation}
取全测试有效网格中 $g$ 最高 10\% 的区域，使用同一 RMSE 公式计分。这个指标不参与主验收，
但用于量化会聚线和尖峰区域的表现。

\section{延迟测试协议}

每个模型先预热 10 次，batch=1。一次计时包括：
\begin{enumerate}
  \item 从一条预构建 NumPy 特征复制并创建 tensor；
  \item 将 tensor 移动到目标设备；
  \item 模型前向；
  \item 残差反归一化并加训练均值场；
  \item 输出复制回 CPU NumPy；
  \item 对 CUDA 执行同步，确保异步 kernel 已完成。
\end{enumerate}
GPU 原始实验统计 200 次；CPU 消融诊断统计 20 次；R1 CPU 独立复核重新统计 200 次。报告
P50、P90、P95 和最大值，验收使用 P95。

% -----------------------------------------------------------------------------
\chapter{五轮消融实验}
\label{chap:ablation}

\section{实验设计}

五轮配置在实验开始前集中登记，共用同一数据、seed、训练预算、验证规则和测试指标：

\begin{longtable}{L{1.2cm}L{3.5cm}L{4.3cm}L{5.2cm}}
\caption{五轮实验假设}\label{tab:ablation-hypothesis}\\
\toprule
轮次 & 主要变化 & 待验证假设 & 实际结论 \\
\midrule
\endfirsthead
\toprule
轮次 & 主要变化 & 待验证假设 & 实际结论 \\
\midrule
\endhead
R1 & hidden 24，modes 12/24 & 小模型能否建立达标基线 & 首轮即达到 0.554 dB，且 CPU 低于 100 ms \\
R2 & hidden 32，modes 16/48 & 更多距离模态能否恢复会聚结构 & RMSE 降至 0.519 dB，但 CPU P95 接近门槛 \\
R3 & 非周期 padding + 残差 & 减少边界环绕与优化困难 & RMSE 最低 0.5165 dB，提升很小，CPU 超限 \\
R4 & Hankel + 3$\times$3 局部分支 & 几何扩散和局部结构能否改善尖峰 & MAE 略降，整体 RMSE 未改善 \\
R5 & 梯度 + 高梯度加权 & 是否能专门压低会聚区误差 & 整体和高梯度 RMSE 均退化 \\
\bottomrule
\end{longtable}

\section{量化结果}

\begin{table}[H]
\centering
\caption{五轮测试精度与复杂度}
\resizebox{\textwidth}{!}{%
\begin{tabular}{lrrrrrrrrr}
\toprule
轮次 & 参数量 & 最佳 epoch & 验证 RMSE & 测试 RMSE & 测试 MAE & 高梯度 RMSE & GPU P95 & CPU P95 & 训练时间 \\
\midrule
R1 & 1.33 M & 180 & 0.5693 & 0.5542 & 0.2163 & 1.5625 & 3.91 ms & 40.06 ms & 99.3 s \\
R2 & 6.30 M & 179 & 0.5437 & 0.5186 & 0.1887 & 1.4771 & 3.80 ms & 95.22 ms & 117.1 s \\
R3 & 6.30 M & 180 & 0.5396 & \textbf{0.5165} & 0.1802 & 1.4799 & 4.73 ms & 143.81 ms & 139.9 s \\
R4 & 6.30 M & 179 & 0.5385 & 0.5169 & \textbf{0.1765} & \textbf{1.4737} & 4.51 ms & 133.17 ms & 168.2 s \\
R5 & 6.30 M & 176 & 0.6506 & 0.6316 & 0.2965 & 1.5184 & 4.50 ms & 131.28 ms & 181.7 s \\
\bottomrule
\end{tabular}}
\end{table}

\figfull{0.98}{fig06_ablation_study.pdf}{五轮精度、时延、复杂度和相对 baseline 改善。}

\section{消融解释}

\subsection{R1：简单模型已经足够}
R1 把学生参考的 64 channels 和 72 modes 大幅缩小，仍取得 0.554 dB。说明任务的关键不是
无上限增加容量，而是标签契约、目标残差化和固定小环境族。

\subsection{R2：模态增加是主要有效改动}
从 12/24 增到 16/48 后，整体 RMSE 降低 0.0356 dB，高梯度 RMSE 降低 0.0854 dB。
这是五轮中最明确的架构收益，但参数量从 1.33 M 增至 6.30 M，CPU P95 升至 95.22 ms。

\subsection{R3：最佳精度但边际收益很小}
padding 和残差把 R2 的 0.51859 dB 降至 0.51652 dB，仅改善约 0.00208 dB。它成为精度
优胜是客观事实，但不意味着必须在所有部署中选择它。

\subsection{R4：物理启发特征不是免费收益}
Hankel 和局部分支把 MAE 从 0.18023 降至 0.17649 dB，高梯度 RMSE略降到 1.47365 dB，
但整体 RMSE 反而略升到 0.51692 dB。对于本固定频率/距离网格，网络与均值场已经能学习
主要几何扩散，附加特征收益不足以抵消复杂度。

\subsection{R5：结构损失是明确负结果}
把高梯度网格权重设为 1.0 后，整体 RMSE 退化到 0.63156 dB，高梯度 RMSE 也退化到
1.51844 dB。可能原因是极窄会聚线存在位置敏感性，强 MSE 权重鼓励模型牺牲大面积平滑区域，
却无法可靠解决峰值错位。该轮证明“更物理的 loss”必须用最终 dB 指标验证，不能只凭直觉保留。

% -----------------------------------------------------------------------------
\chapter{最终测试结果与独立验收}
\label{chap:results}

\section{R3 精度优胜模型}

在 64 条测试样本、1,567,951 个有效网格上：

\begin{table}[H]
\centering
\caption{R3 测试集详细指标}
\begin{tabularx}{\textwidth}{L{5.0cm}C{3.0cm}Y}
\toprule
指标 & 数值 & 解释 \\
\midrule
整体 RMSE & 0.5165 dB & 主精度指标 \\
整体 MAE & 0.1802 dB & 平均绝对偏差 \\
Bias & -0.0008 dB & 几乎无系统性正负偏差 \\
绝对误差 P50 / P90 / P95 & 0.0679 / 0.3510 / 0.6834 dB & 95\% 网格低于 0.684 dB \\
绝对误差 P99 / P99.9 & 2.0611 / 6.4881 dB & 尾部来自窄会聚线 \\
逐样本 RMSE 中位 / P90 & 0.5120 / 0.6865 dB & 样本级稳定性 \\
最差单样本 RMSE & 0.7687 dB & 仍远低于 2 dB \\
高梯度区域 RMSE & 1.4799 dB & 困难网格仍通过 2 dB 参考线 \\
误差 $>2$ / $>5$ / $>10$ dB 网格 & 1.049\% / 0.206\% / 0.018\% & 大误差空间占比很小 \\
\bottomrule
\end{tabularx}
\end{table}

\section{与 baseline 的分布级比较}

R3 相比训练均值场 baseline 将整体 RMSE 从 1.4211 降至 0.5165 dB，降幅 63.65\%。
更关键的是，模型在每条测试样本上都显著低于 baseline，改善不是由少数简单样本拉低平均值。

\figfull{0.98}{fig07_error_distribution.pdf}{R3 与训练均值场 baseline 的网格误差 CDF 和逐样本 RMSE。}

\section{空间误差结构}

平均误差图呈现清晰的会聚线结构，说明主要误差来自预测会聚边界相对 Bellhop 的轻微位置
偏移，而不是整个区域的均匀偏置。误差随距离整体上升，符合 50 km 远距离结构更复杂的直觉；
但平均误差仍保持较低水平。

\figfull{0.98}{fig08_error_spatial_analysis.pdf}{测试集逐网格平均/P95 误差及其距离、深度分解。}

\section{最差测试样本}

最差样本为 \code{ssp\_00154}，四参数依次为 -1.4514 m/s、-1.7405 m/s、+31.63 m、
+0.4432 m/s。其逐场 RMSE 为 0.7687 dB。最差有效单点位于约 46.48 km、1781.6 m：
Bellhop 为 68.30 dB，模型为 87.75 dB，绝对误差 19.44 dB。该点位于非常窄的会聚线，
单点误差很大，但周围大多数网格仍准确，因此整场 RMSE 低于 0.8 dB。

\figfull{0.98}{fig09_worst_case.pdf}{最差测试样本的参考场、预测场、误差图、切片和误差直方图。}

\section{独立新进程复核}

为避免只信任训练进程内的数字，项目在新 Python 进程中执行以下步骤：重新计算数据 SHA、
重新加载 checkpoint、重建特征、只读取冻结 test、重新推理并重新计时。

\begin{table}[H]
\centering
\caption{独立复核结果}
\begin{tabularx}{\textwidth}{L{4.4cm}C{3.0cm}C{3.0cm}Y}
\toprule
复核项 & 原记录 & 新进程 & 结论 \\
\midrule
R3 CUDA RMSE & 0.516517 & 0.516512 & 差 4.91$\times10^{-6}$ dB \\
R3 CUDA P95 & 4.73 ms & 4.27 ms & 通过 \\
数据 SHA & \multicolumn{2}{c}{完全一致} & 通过 \\
R1 CPU RMSE & CUDA 0.5542 & CPU 0.7582 & 跨后端重新验收通过 \\
R1 CPU MAE & CUDA 0.2163 & CPU 0.3067 & 通过 \\
R1 CPU P95 & 训练诊断 40.06 ms & 200 次复核 38.32 ms & 通过 \\
\bottomrule
\end{tabularx}
\end{table}

R1 CPU 输出与 CUDA 输出存在 0.204 dB RMSE 差异，可能来自多层 complex FFT 与 einsum 后端
以及 CUDA TF32 数值路径差异。因此 CPU 复核按“同一权重在目标后端重新验收”处理，不宣称
与 CUDA 逐值一致。CPU 自身的精度和时延仍有充足余量。

% -----------------------------------------------------------------------------
\chapter{性能、部署与工程交付}
\label{chap:deployment}

\section{观测计算加速}

正式 Bellhop 单场中位时间 9.138 s。以完整输出同一 96$\times$256 TL 场计算：
\begin{align}
  S_{\mathrm{R3,A100}} &\approx \frac{9138\ \mathrm{ms}}{4.732\ \mathrm{ms}}\approx 1931,\\
  S_{\mathrm{R1,CPU}} &\approx \frac{9138\ \mathrm{ms}}{38.317\ \mathrm{ms}}\approx 238.
\end{align}
这是实际墙钟观测比，不是同硬件同实现的纯算法复杂度比；Bellhop 和网络运行后端不同。
即便如此，它直接反映了在线业务所能获得的数量级加速。

\section{部署选择}

\noindent
\begin{tabularx}{\textwidth}{L{4.0cm}Y Y}
\toprule
场景 & 推荐 & 理由 \\
\midrule
甲方已有 CUDA GPU & R3 & 最低 RMSE，P95 约 4--5 ms \\
仅 CPU 或硬件未冻结 & R1 & 11 MB、1.33 M 参数、独立 CPU P95 38.32 ms \\
极端重视高梯度 MAE & 不建议单独选择 R4 & 收益很小且 CPU 超限 \\
希望继续加入结构 loss & 暂停 & R5 已给出明确负结果 \\
\bottomrule
\end{tabularx}

\section{模型输入输出契约}

上线服务应固定：
\begin{itemize}
  \item 输入 SSP 深度范围覆盖 0--2000 m，并插值到冻结 96 点接收深度；
  \item 输入单位必须是 m 和 m/s；
  \item 输出为非相干 TL，单位 dB，形状 96$\times$256；
  \item 距离坐标固定 0.1--50 km，深度坐标固定 10--1990 m；
  \item 使用 checkpoint 内保存的训练均值场和残差尺度；
  \item 不在服务端重新估计归一化统计。
\end{itemize}

\section{版本与哈希}

\noindent
\begin{tabularx}{\textwidth}{L{4.1cm}Y}
\toprule
产物 & 版本/哈希 \\
\midrule
代码仓库 & \url{https://github.com/xuangu-fang/ocean-acoustic-surrogate}（private） \\
最终 MVP tag & \code{mvp-v0.1} \\
数据契约 tag & \code{dataset-contract-v0.1} \\
数据 profile tag & \code{dataset-v0.1} \\
数据 SHA-256 & \footnotesize\nolinkurl{13f4246b86a2fff33f2e5d587c1443be47eb15b4ad3d5c413688464ed8c7281d} \\
R1 checkpoint SHA-256 & \footnotesize\nolinkurl{fca7a048d3611c8587d7d8395ef38332312ee1c489d101052603e640fefa5208} \\
R3 checkpoint SHA-256 & \footnotesize\nolinkurl{83c8b233a8d272d4aed1a50267b71b00afee2e6497a8988e6f65b2f4e9449048} \\
\bottomrule
\end{tabularx}

% -----------------------------------------------------------------------------
\chapter{风险、适用边界与后续建议}
\label{chap:risks}

\section{当前成果可以支持的结论}

\begin{passbox}[可以确认]
在冻结的 1 kHz、50 km、2000 m 水深、50 m 源深、固定沙质海底和窄 SSP 环境族内，
高质量 Bellhop 非相干标签可以用 512 样本训练出显著满足 2 dB / 100 ms 的 FNO 代理；
结果已通过密封测试和独立新进程复核。
\end{passbox}

\section{当前成果不能外推的结论}

\begin{riskbox}[不可直接外推]
本文没有证明跨频率、跨源深、跨水深、跨底质、跨海区、距离相关环境、相干/半相干场或
甲方私有 Bellhop 设置上的泛化能力。超出冻结参数的输入应被服务层拒绝或明确标记为域外，
而不是默默返回一个看似合理的场。
\end{riskbox}

\section{甲方 Bellhop 对齐风险}

甲方 Bellhop 的射线数、射线角、波束类型、插值、场模式、TL 下限和无效网格处理目前未知。
即使环境输入相同，也可能产生系统性差异。建议在大规模补数据前执行最小对齐：

\begin{enumerate}
  \item 从本数据族选择 8--16 条覆盖四参数边界和中心的 SSP；
  \item 双方使用完全相同的环境、网格和输出模式；
  \item 比较整体 RMSE、bias、P95、会聚位置和无效网格；
  \item 若主要是全局偏差，优先做输出校准；
  \item 若是场模式或射线离散差异，先统一数值契约，再考虑微调模型。
\end{enumerate}

\section{推荐的最小后续路线}

\begin{enumerate}
  \item \textbf{阶段 A：}甲方 Bellhop 同 case 对齐，不扩展数据维度；
  \item \textbf{阶段 B：}若必要，以 64--256 条甲方标签做校准/微调；
  \item \textbf{阶段 C：}只有在首版交付稳定后，才逐一加入源深、频率或距离相关环境；
  \item \textbf{阶段 D：}每增加一个维度，都保留原窄域测试集作为回归门槛，防止泛化扩展
  损害已达标场景。
\end{enumerate}

% -----------------------------------------------------------------------------
\chapter{结论}

本项目用一个刻意收窄但严格满足硬参数的深海环境族，完成了可复现、可审计、可交付的代理
模型闭环。数据侧不是简单宣称“Bellhop 即真值”，而是先完成射线数收敛；模型侧不是盲目
堆叠更大 FNO，而是用五轮消融确认哪些改动有收益、哪些是负结果；测试侧同时报告 RMSE、
MAE、尾部、逐样本、高梯度和 P95 延迟，并用新进程重载验证结果。

最终结果给出清晰选择：R3 在 A100 上提供约 0.52 dB 的最佳精度，R1 在普通 CPU 上仍以
0.76 dB / 38 ms 大幅通过验收。当前最合理的工程决策不是继续增加数据覆盖或结构复杂度，
而是先与甲方 Bellhop 完成少量同 case 对齐，并把已达标窄域模型作为稳定基线。

% -----------------------------------------------------------------------------
\appendix
\chapter{复现命令}

\begin{lstlisting}[style=command]
cd /home/ubuntu/project/ocean-acoustic-surrogate
uv sync --locked --group dev
uv run pytest

export OCEAN_SURROGATE_ROOT=/home/ubuntu/ocean-acoustic-surrogate-artifacts

# 标签收敛、正式数据和 profile
uv run ocean-acoustic-surrogate pilot configs/mvp.yaml --samples 8
uv run ocean-acoustic-surrogate generate configs/mvp.yaml --samples 512
uv run ocean-acoustic-surrogate profile configs/mvp.yaml --samples 512

# 五轮消融
uv run ocean-acoustic-surrogate campaign \
  configs/mvp.yaml configs/campaign.yaml --samples 512

# R3 CUDA 独立复核
uv run ocean-acoustic-surrogate verify configs/mvp.yaml \
  /home/ubuntu/ocean-acoustic-surrogate-artifacts/runs/\
  20260827T104438Z_r3_padding_residual \
  --samples 512 --device cuda

# R1 CPU 独立复核
uv run ocean-acoustic-surrogate verify configs/mvp.yaml \
  /home/ubuntu/ocean-acoustic-surrogate-artifacts/runs/\
  20260827T104036Z_r1_fno_small \
  --samples 512 --device cpu

# 重建本白皮书
bash scripts/build_whitepaper.sh
\end{lstlisting}

\chapter{产物目录}

\noindent
\begin{tabularx}{\textwidth}{L{5.0cm}Y}
\toprule
内容 & 路径 \\
\midrule
代码与配置 & \path{/home/ubuntu/project/ocean-acoustic-surrogate} \\
完整数据 & \path{/home/ubuntu/ocean-acoustic-surrogate-artifacts/datasets/scs_june_narrow_v0.1/n512} \\
R1 完整 run & \path{/home/ubuntu/ocean-acoustic-surrogate-artifacts/runs/20260827T104036Z_r1_fno_small} \\
R3 完整 run & \path{/home/ubuntu/ocean-acoustic-surrogate-artifacts/runs/20260827T104438Z_r3_padding_residual} \\
机器可读结果 & \path{docs/results} \\
白皮书图表 & \path{docs/whitepaper/assets} \\
白皮书 PDF & \path{docs/Ocean_Acoustic_Surrogate_Whitepaper_v1.0.pdf} \\
\bottomrule
\end{tabularx}

\chapter{术语表}

\begin{longtable}{L{3.2cm}L{11.0cm}}
\toprule
术语 & 含义 \\
\midrule
\endfirsthead
\toprule
术语 & 含义 \\
\midrule
\endhead
SSP & Sound Speed Profile，声速随深度的剖面 \\
TL & Transmission Loss，传播损失，单位 dB \\
Bellhop & 基于射线/高斯束的水声传播数值模型 \\
非相干 TL & 路径强度相加而不保留路径间相位干涉的传播损失 \\
FNO & Fourier Neural Operator，Fourier 神经算子 \\
RMSE & 均方根误差，对较大误差更敏感 \\
MAE & 平均绝对误差，直观反映平均偏差 \\
P95 & 95\% 观测不超过的分位数 \\
密封测试集 & 在训练和 checkpoint 选择中不使用标签的固定测试划分 \\
高梯度区域 & 参考 TL 深度/距离有限差分幅度最高 10\% 的有效网格 \\
残差学习 & 预测相对训练均值场的变化，而非直接预测完整 TL \\
\bottomrule
\end{longtable}

\chapter{参考资料}

\begin{enumerate}
  \item Z. Li et al., ``Fourier Neural Operator for Parametric Partial Differential Equations,'' ICLR, 2021.
  \item M. B. Porter, \textit{The BELLHOP Manual and User's Guide: Preliminary Draft}, 2011.
  \item F. B. Jensen, W. A. Kuperman, M. B. Porter, and H. Schmidt, \textit{Computational Ocean Acoustics}, 2nd ed., Springer, 2011.
\end{enumerate}

\vfill
\begin{center}
\color{gray}\small
—— 文档结束 ——\\[0.5em]
所有指标均来自冻结产物；详细 JSON、checkpoint 哈希和构建脚本随项目仓库交付。
\end{center}

\end{document}
